\chapter{연구 내용}
\section{모방학습 실험과 CV+IK 전환}
\subsection{단일 블록 시연의 범위 축소와 재확장}
초기에는 리더암의 시연으로 블록 조작 동작을 학습하도록 설계하였다. ACT(Action Chunking with Transformers)는 연속된 여러 시점의 행동을 묶어 예측하는 모방학습 방법이며\cite{act}, SmolVLA는 영상·언어 입력을 행동 생성에 연결하는 모델이다\cite{smolvla}. LeRobot으로 시연 수집과 학습·추론 환경을 구성하고, 로봇 제어기와 정책 서버는 gRPC 원격 추론으로 연결하였다.

첫 실험은 작업 영역 전체에 단일 블록을 무작위로 놓고, 이를 텔레오퍼레이션으로 집어 중앙 목표 지점에 옮기는 것이었다. 약 100개 에피소드를 수집하여 ACT와 SmolVLA를 학습·실행했으나, 팀이 관찰한 수행 성공률은 약 2--10\%로 낮았다. 따라서 처음부터 넓은 시작 위치를 다루기보다 한정된 조건에서 동작을 학습할 수 있는지 먼저 확인하였다.

두 번째 단계에서는 특정 시작 위치에서 목표 지점까지 옮기는 시연을 중심으로 51개 에피소드를 수집하였다. 이 조건에서 단일 블록 이동은 50회 중 29회(58\%) 성공하였다. 위치를 한정하면 수행이 나아지는 것을 보고 시작 위치를 다시 전 영역으로 넓혔으며, 약 120개 에피소드로 ACT와 SmolVLA를 학습하였다. 그러나 이때의 수행 성공률은 약 10--20\%로 다시 낮아졌다. 세 단계 모두 블록 하나를 집어 옮기는 과제로, 블록 수를 늘리지 않고 시작 위치의 범위를 바꾸었다.

\begin{table}[htbp]
\centering
\caption{초기 단일 블록 모방학습의 조건과 수행 결과}\label{tab:early}
\begin{tabularx}{\linewidth}{@{}L{0.14\linewidth}YL{0.17\linewidth}L{0.23\linewidth}@{}}
\toprule 단계 & 블록의 시작 조건 & 시연 에피소드 & 수행 성공 \\
\midrule
전 영역 수집 & 단일 블록, 전 영역 무작위 위치 & 약 100개 & 약 2--10\% \\
범위 축소 & 단일 블록, 특정 시작 위치 중심 & 51개 & 29/50회 (58\%) \\
범위 재확장 & 단일 블록, 전 영역 무작위 위치 & 약 120개 & 약 10--20\% \\
\bottomrule
\end{tabularx}
\end{table}

표~\ref{tab:early}의 51개 시연과 29/50회는 실험 집계값이며, 나머지 수집량과 성공률 범위는 팀의 회고에 따른 근사치이다. 단계별 작업 범위와 학습 조건이 달라 데이터 수 증가만의 효과로 해석하지 않았으며, 모델별 세부 성공률을 구분할 기록도 확보하지 못하였다.

\subsection{ACT와 CV+IK의 하이브리드 시도}
단일 블록의 파지와 이동에서도 성공이 불안정하자, 학습 정책이 전체 동작을 담당하는 부담을 줄이기 위해 ACT와 CV+IK의 하이브리드 방식을 시도하였다. 파지에는 ACT를 사용하고 배치는 IK로 연결하는 역할 분담을 검토·시험하였다. SmolVLA 대신 ACT를 선택한 이유는 팀의 비교 실험에서 추론이 더 빠르고 수행 성공률도 일관되게 높았기 때문이다. 이는 당시 장비와 학습 조건에서의 관찰이며, 모델별 지연시간과 성공률을 같은 조건에서 측정한 비교표는 확보하지 못하였다.

ACT의 학습·추론 설정을 조정했지만 약 60\%를 넘는 안정적인 수행으로 끌어올리기 어려웠다. 팔이 목표 근처에 도달한 뒤에도 블록을 잡지 못하는 경우에는 영상 입력, 정책 출력과 실제 접촉 중 어느 부분을 바꿔야 할지 판단하기 어려웠다. 하이브리드로 정책의 역할을 줄여도 파지의 불안정성이 남으면 반복 미션 전체를 안정화하기 어려웠다.

\subsection{CV+IK 기반 제어로의 전환}
카메라를 조금 옮긴 뒤 기존 정책의 수행이 달라져 시연부터 다시 수집해야 했으며, 100개 이상의 에피소드를 사람이 다시 녹화하는 부담이 컸다. 장비를 시연장으로 옮긴 뒤에도 같은 수집·학습·평가 과정을 반복해야 할 수 있어 현장에서 조정할 수 있는 제어 방식이 필요하였다.

이에 크기와 색상이 정해진 블록과 고정 작업 평면에 맞추어 CV+IK 제어를 구현하였다. 반복적인 수동 시연을 줄이기 위해 실행 과정도 학습용으로 기록하고자 하였으며, 동작을 보완한 뒤 이를 Task3의 자동 수집 기능으로 구현하였다. 저장한 자료를 이용한 정책 재학습은 후속 과제로 남았다.

CV는 블록 중심과 방향을, 캘리브레이션은 로봇 좌표를, IK는 후보 자세와 도달성을 제공한다. 그리퍼 위치·부하는 파지 여부를 확인하고 FSM은 재시도와 대상 전환을 관리한다. 중간 결과를 분리하면 검출 실패, 좌표 편차, 계획 불가와 빈손 파지를 각각 점검할 수 있다. 기존 모터 통신과 장비 제어 경험은 이 경로에 재사용하였다.

\begin{table}[htbp]
\centering
\caption{초기 실험에서 관찰한 제약과 전환 후의 설계 대응}
\begin{tabularx}{\linewidth}{@{}YY@{}}
\toprule 초기 운용에서의 제약 & CV+IK의 대응 \\
\midrule
시작 위치 확대 후 단일 블록 수행 불안정 & 검출 위치를 로봇 좌표로 변환하고 후보별 IK 도달성 확인 \\
카메라 변경 후 반복 시연 수집 부담 & 카메라 재고정·재캘리브레이션 후 실행 자료를 Task3으로 기록 \\
실패 원인의 분리 어려움 & 검출·좌표·IK·센서·상태 전이의 중간 결과 기록 \\
하이브리드에서도 남은 파지 실패 & 센서 파지 확인 뒤 운반, 회전 재시도와 대상 재선택 \\
\bottomrule
\end{tabularx}
\end{table}

CV+IK에서도 카메라 고정과 마운트 변경 후 재캘리브레이션, 조명에 따른 검출 조정이 필요하다. 환경 변화의 영향을 없앤다는 가정 대신, 관측과 제어를 점검·보정하고 실행 자료를 다시 수집할 수 있는 구조를 선택하였다. 두 방법의 일반화 성능을 동일 조건에서 비교하는 실험은 수행하지 않았다.
\FloatBarrier

\section{최종 시스템 아키텍처}
\subsection{구성 요소와 관측 경로}
\label{sec:system-config}
기본 제어는 고정 상단 카메라의 블록 검출, 평면 좌표 변환과 IK, 그리퍼 위치·부하 판정으로 구성된다. 미션 평가에는 ArUco 위치 추정과 손목 카메라 색상 게이트를 추가 관측 기능으로 사용하였다. 이하에서는 기본 제어 구조와 추가 관측 기능의 역할을 구분하여 설명한다.

시스템은 카메라 서비스, 인식 모듈, 기구학·동작 모듈, 상태 제어기와 로봇 입출력으로 구성된다. 카메라 서비스가 USB 장치를 단독으로 읽어 최신 프레임을 제공하고, 미션 실행기는 HTTP로 받은 최신 정지 영상(snapshot)에서 제어할 블록을 검출한다. 브라우저 관제 화면에는 MJPEG 영상과 분석 결과를 표시한다. 화면에 겹쳐 그린 좌표는 운영자의 점검용이며 파지 계획의 입력으로 사용하지 않는다.

\begin{figure}[htbp]
\centering
\begin{tikzpicture}[node distance=9mm,
box/.style={draw,rounded corners,align=center,text width=5.4cm,minimum height=9mm,font=\small},
arr/.style={-{Latex},thick}]
\node[box] (cam) {고정 카메라 / 카메라 서비스};
\node[box,below=of cam] (cv) {최신 snapshot / 고전 CV\\영역 밖 블록 중심·방향·색};
\node[box,below=of cv] (ik) {베이스 좌표 / 파지 후보 생성\\IK 도달성 검사와 웨이포인트};
\node[box,below=of ik] (fsm) {FSM / 관절 궤적 실행\\파지 확인·운반·배치·재시도};
\node[box,below=of fsm] (robot) {SO-101 팔로워\\관절 위치·그리퍼 부하};
\node[box,text width=3.4cm,right=13mm of cv] (ui) {브라우저 관제\\MJPEG·검출·FK 표시};
\draw[arr] (cam)--(cv);
\draw[arr] (cv)--(ik);
\draw[arr] (ik)--(fsm);
\draw[arr] (fsm)--(robot);
\draw[arr] (cam.east)-|(ui.north);
\draw[arr,dashed] (robot.west)--++(-9mm,0)|-node[pos=0.75,left,font=\scriptsize]{측정 피드백}(fsm.west);
\draw[arr,dashed] (robot.east)-|(ui.south);
\end{tikzpicture}
\caption{기본 CV+IK 제어 경로와 관제 구성. 실선은 주 입력·실행 경로, 점선은 측정 상태의 전달이다. 추가 관측 기능은 절~\ref{sec:additional-observation}에서 설명한다.}
\label{fig:architecture}
\end{figure}

\FloatBarrier
\subsection{전체 상태 흐름과 판단 조건}
제어 상태는 SELECT, PICK, VERIFY, TRANSPORT, PLACE 순으로 이어진다. SELECT는 카메라 시야를 확보하는 대기 자세(HOME)에서 최신 영상을 받아 블록을 고른다. PICK은 도달 가능한 파지 후보를 실행하고 센서값을 확인한다. PICK이 통과해도 VERIFY에서 파지 상태를 다시 검사하며, 물체가 확인된 경우에만 TRANSPORT로 넘어간다. 실패하면 후보를 바꾸거나 SELECT로 돌아간다(그림~\ref{fig:fsm}).

\begin{figure}[htbp]
\centering
\begin{tikzpicture}[
state/.style={draw,rounded corners,text width=3.1cm,minimum height=10mm,align=center,font=\small},
decision/.style={draw,diamond,aspect=2.5,text width=2.1cm,align=center,inner sep=1pt,font=\small},
arr/.style={-{Latex},thick}]
\node[state] (sel) at (0,0) {SELECT\\관측과 대상 선택};
\node[state] (pick) at (0,-1.7) {PICK\\후보별 파지 시도};
\node[decision] (ver) at (0,-3.4) {VERIFY\\파지 확인?};
\node[state] (tr) at (0,-5.1) {TRANSPORT\\슬롯 상공 이동};
\node[state] (pl) at (0,-6.5) {PLACE\\블록 해제};
\node[state,text width=2.5cm] (retry) at (-4.3,-3.4) {재시도 / 순회 보류\\시도 이력 관리};
\node[state,text width=2.7cm] (done) at (5.4,0) {DONE\\관측 기반 완료};
\draw[arr] (sel)--node[right,font=\scriptsize]{대상 있음}(pick);
\draw[arr] (pick)--(ver);
\draw[arr] (ver)--node[right,font=\scriptsize]{예}(tr);
\draw[arr] (tr)--(pl);
\draw[arr] (ver.west)--node[above,font=\scriptsize]{아니오}(retry.east);
\draw[arr] (pick.west)-|node[pos=0.25,above,font=\scriptsize]{PICK 실패}(retry.north);
\draw[arr] (retry.west)--++(-5mm,0)|-(sel.west);
\draw[arr] (pl.east)--++(8mm,0)|-(sel.south east);
\draw[arr] (sel)--node[above,align=center,font=\scriptsize]{영역 밖 검출\\연속 5초 부재}(done);
\draw[arr] (sel.north east)--++(0,6mm)-|(sel.north west);
\node[font=\scriptsize,anchor=south] at (0,1.15) {대상 미확정 / 새 프레임 재관측};
\end{tikzpicture}
\caption{1차 미션의 상태 흐름과 성공/실패 분기. 완료 판정에는 유효한 새 프레임만 사용한다. 외부 시간 제한에 의한 중단은 완료와 구분한다.}
\label{fig:fsm}
\end{figure}

\FloatBarrier

\subsection{외부 기술과 팀의 구현 범위}
\begin{table}[htbp]
\centering
\caption{외부 기술의 활용과 팀의 설계 및 구현}
\begin{tabularx}{\linewidth}{@{}L{0.16\linewidth}L{0.29\linewidth}Y@{}}
\toprule
구성 & 활용한 기반 기능 & 팀의 설계 및 구현 \\
\midrule
인식 & OpenCV 색 변환, 윤곽선과 평면 변환 & 미터법 형상 필터, 색 기준점 할당, 영역 밖 대상 선택 \\
기구학 & PlaCo의 URDF 기반 FK/IK 풀이 & 시드와 중립 손목 방향, 후보 보정, 접근 높이 탐색 \\
장비 제어 & LeRobot 장비 통신과 관절 명령 제한 & 입출력 추상화, 측정 자세 기반 보간, 파지 판정 연결 \\
미션 실행 & Python 상태 및 자료 구조 & FSM, 후보 재시도와 순회 보류, 동적 슬롯, 적재 PLACE 전략 \\
관제·기록 & 영상 입출력과 HTTP 기본 기능 & 카메라 단독 소유, 최신 프레임 전달, FK 표시, 실험 이력 관리 \\
\bottomrule
\end{tabularx}
\end{table}

카메라 장치의 접근 충돌을 막기 위해 영상 획득을 별도 서비스로 분리하였다. 관제 화면의 FK 표시는 제어기가 공유한 관절 측정값을 사용하며, 영상 위에 검출 후보와 FK 투영 위치를 함께 표시한다.

\section{블록 검출과 대상 선택}
\subsection{미터법 영상과 후보 생성}
지정 영역의 빨간 테이프가 빨간 블록과 함께 검출되는 문제가 있었다. 색 임계값만 사용하면 같은 색의 경계가 후보에 남는다. 블록의 채워진 정사각형 윤곽과 길쭉하거나 내부가 빈 테이프 윤곽을 구분하도록 형상 필터를 추가하였다. 색상으로 후보를 찾은 뒤 기하 특성으로 거르는 방식이다.

입력 영상은 캘리브레이션 결과에 따라 로봇 평면 좌표에 대응하는 1\,mm/pixel 정사영 영상으로 변환한다. HSV 색 공간에서 색별 마스크를 만든 뒤 타원형 커널로 열림·닫힘 연산을 적용한다. 외부 윤곽선에서 면적, 회전 사각형과 형상 지표를 계산하므로 물체 크기 조건을 밀리미터 단위로 지정할 수 있다.

검출기의 주요 설정값은 표~\ref{tab:detector-config}와 같다. 단일 블록의 이상적인 상면 면적은 $40\times40=1,600\,\mathrm{mm}^2$이며, 측면·그림자·마스크 손실을 고려하여 후보 면적 범위에 여유를 두었다.

\begin{table}[htbp]
\centering
\caption{검출기의 설정값과 역할}\label{tab:detector-config}
\begin{tabularx}{\linewidth}{@{}L{0.31\linewidth}L{0.19\linewidth}Y@{}}
\toprule
항목 & 설정값 & 적용 목적 \\
\midrule
정사영 해상도 & 1\,mm/pixel & 형상 조건을 미터법으로 표현 \\
후보 면적 & 900--3,200\,$\mathrm{mm}^2$ & 작은 잡음·큰 배경 영역 배제 \\
최대 종횡비 & 1.6 & 길쭉한 테이프 후보 배제 \\
최소 채움 비율 & 0.65 & 비어 있는 경계 형태 배제 \\
최소 볼록도 & 0.78 & 지나치게 오목한 형상 배제 \\
형태학적 커널 & $9\times9$ pixel & 마스크 잡음과 윤곽 변동 완화 \\
중복 후보 거리 & 30\,mm 미만 & 겹친 색 마스크의 동일 물체 병합 \\
색상별 최대 후보 & 1개 & 과제의 색상별 단일 블록 조건 반영 \\
관측 부채꼴 & 반경 320\,mm, $-90^\circ$--$90^\circ$ & 제어 대상 관측 범위 제한 \\
\bottomrule
\end{tabularx}
\end{table}

후보는 면적, 종횡비와 채움 비율로 판별한다. 맞닿은 블록의 윤곽이 합쳐지거나 팔에 가려져 면적이 줄면 실제 블록도 제외될 수 있다. 이를 줄이기 위해 SELECT에서는 HOME으로 돌아가 시야를 확보하고 매 사이클 새 영상을 검출한다. 형상 필터 적용 전후의 검출 정확도는 따로 측정하지 않았다.

\subsection{색 기준점 할당과 영역 필터}
후보의 대표 색상·채도를 $(H,S)$, 기준점을 $(H_c,S_c)$라고 하면 구현의 거리 함수는 다음과 같다.
\begin{equation}
\Delta H=\min(|H-H_c|,180-|H-H_c|),\qquad
d=\sqrt{\left(\frac{2\Delta H}{128}\right)^2+
\left(\frac{S-S_c}{128}\right)^2}.
\end{equation}
색상은 원형 범위로 처리하고, 한 색에 기준점이 여러 개 있으면 가장 가까운 거리를 사용한다. 전체 후보--색 쌍을 거리순으로 정렬한 후 이미 할당한 후보나 사용한 색을 제외하면서 탐욕적으로 할당한다. 거리가 0.35를 넘는 후보는 미할당 상태로 남긴다.

동일한 블록이 여러 색 마스크에 잡힌 경우에는 채움 비율과 볼록도의 곱이 높은 후보를 우선하여 가까운 후보를 병합한다. 지정 영역 내부 후보는 블록 중심 좌표로 판정하여 색 할당 전에 제외한다. 이 순서를 통해 이미 배치한 블록이 색 할당을 차지하여 영역 밖의 후보를 가리는 상황을 줄였다.

1차 미션은 현재 검출된 후보 가운데 베이스 원점에서 가까운 블록을 먼저 선택한다. 색상별 블록이 하나씩이라는 조건을 이용하여 색 이름을 재시도 식별자와 슬롯 예약의 키로 사용한다. 동일한 색이 여러 개 존재하는 일반 작업으로 확장하려면 별도의 물체 추적 식별자가 필요하다.

\section{캘리브레이션과 좌표 관리}
\subsection{대응점 수집 절차}
영상의 블록 중심과 실제 그리퍼 접근 위치가 어긋나지 않도록 대응점 수집 조건을 정하였다. 블록 윗면과 빈 테이블은 높이가 다르며, 같은 위치에서 턱이 닿아도 손목 회전에 따라 FK 기준점 좌표가 달라질 수 있다. 대응점 수뿐 아니라 기준 평면과 손목 자세, 기록 항목을 함께 관리하였다.

대응점 하나를 기록할 때 같은 블록 위치에서 두 장면을 촬영하였다. 먼저 그리퍼 기준점을 블록 윗면 중앙에 맞추고 관절값과 접촉 자세 영상을 저장한다. 블록을 그대로 둔 채 팔을 물린 뒤 윗면이 드러난 영상에서 픽셀 중심을 찾는다. 첫 장면의 관절값으로 구한 FK 좌표와 두 번째 장면의 영상 좌표를 같은 점으로 연결한다. 그림~\ref{fig:calibration-pair}는 초기 수집 자료의 대응점 P4로 절차를 보여준다. 최종 보정에 사용한 대응점은 카메라 재고정 후 별도로 수집하였다.

\begin{figure}[!htbp]
\centering
\begin{subfigure}[t]{0.48\linewidth}
\centering
\includegraphics[width=\linewidth,trim=400bp 200bp 480bp 220bp,clip]{calibration_p4_contact.jpg}
\caption{접촉 자세에서 관절값 기록}
\small FK $(x,y)=(234.237,-43.885)$\,mm
\end{subfigure}\hfill
\begin{subfigure}[t]{0.48\linewidth}
\centering
\begin{tikzpicture}
\node[anchor=south west,inner sep=0] (photo) at (0,0) {\includegraphics[width=\linewidth,trim=400bp 200bp 480bp 220bp,clip]{calibration_p4_clear.jpg}};
\begin{scope}[x={(photo.south east)},y={(photo.north west)}]
\draw[cyan!80!blue,line width=0.9pt] (0.66954,0.46207) circle[radius=2mm];
\draw[cyan!80!blue,line width=0.9pt] (0.645,0.46207)--(0.695,0.46207) (0.66954,0.43)--(0.66954,0.495);
\end{scope}
\end{tikzpicture}
\caption{팔을 물린 뒤 픽셀 중심 확인}
\small $(u,v)=(667.817,381.379)$\,pixel
\end{subfigure}
\caption{대응점 수집 절차의 예. 접촉 자세의 관절값으로 FK 좌표를 구하고, 팔을 물린 영상에서 같은 블록의 픽셀 중심을 기록한다. 두 사진은 원본의 같은 범위를 확대한 것이다.}
\label{fig:calibration-pair}
\end{figure}

수집 도구는 점 이름, 사진, 픽셀 좌표, FK 좌표와 관절값을 함께 저장하며 중단된 지점부터 수집을 재개할 수 있다. 베이스 좌표는 밀리미터로 관리하고 IK 라이브러리 경계에서 단위를 변환한다. 카메라 웹 뷰어의 검출 후보와 FK 투영점은 대응 관계를 점검하는 표시일 뿐, 실제 턱 위치로 사용하지 않는다.

손목 회전은 중립 근처로 유지한다. 같은 턱 접촉 상태라도 회전축과 URDF 그리퍼 기준점의 위치 관계 때문에 기록되는 FK 좌표가 바뀔 수 있기 때문이다. 따라서 좌표값만 남기는 대신 관절각을 보관하고, 캘리브레이션 자세와 런타임 파지 자세의 차이를 줄이는 절차를 사용하였다. 카메라를 다시 고정한 경우에는 대응점을 다시 수집하여 변경된 카메라 배치에 맞는 변환을 적용하였다.

\subsection{좌표 품질과 카메라 안정성 평가}
좌표 변환은 적합 RMS와 최대 LOO 오차로 평가하였다. 카메라 배치나 수집 자세가 바뀔 때 대응점을 재수집하였으며, 최종 평가에는 아홉 점으로 적합한 변환을 사용하였다. 좌표 정합 결과는 절~\ref{sec:calibration-results}에 제시하며, 실제 파지 여부는 그리퍼 센서와 미션 수행 결과로 확인하였다.

카메라 안정성은 기준 영상에 대한 대응 코너의 이동량으로 점검한다. 일부 코너의 검출 변동에 민감한 최댓값 대신 전체 변위 분포의 p95를 사용하였다. 카메라 고정 상태가 달라지면 기존 좌표 변환의 적용 여부도 다시 확인해야 한다.
\FloatBarrier
\section{역기구학과 궤적 실행}

\subsection{손목 과열에 대응한 중립 방향 선택}
초기에는 그리퍼 yaw를 베이스 기준으로 고정하였다. 팔이 선회할수록 손목이 반대로 비틀렸고, 결국 손목 서보가 과열 보호로 정지하였다. 이후에는 작업 위치별 중립 손목 방향을 기준으로 블록 면을 정렬하여 파지 방향을 유지하면서 손목 회전을 줄였다.

파지 계획은 선택한 블록의 좌표와 방향, 캘리브레이션의 파지 높이를 입력으로 받는다. URDF에 정의된 그리퍼 기준 좌표계를 목표 프레임으로 사용하고 접근축이 아래를 향하는 자세를 우선한다. 위치마다 가능한 관절 조합이 다르므로, FK로 구성한 자세 시드를 사용하고 어깨 선회각에 여러 오프셋 후보를 적용하여 해를 찾는다. 선회각의 오프셋 후보는 $0^\circ,\pm6^\circ,\pm12^\circ$이다.

블록 방향을 지정하지 않은 IK는 해당 위치의 중립 손목 방향을 기준으로 한다. 블록 면에 맞춰 잡을 때는 검출 윤곽의 최소 면적 회전 사각형에서 블록 방향 $\theta_b$를 구한다. 최종 파지 방향은 베이스 중심 작업 영역의 접선 방향 $\theta_t=\operatorname{atan2}(y,x)+90^\circ$에 가장 가까운 정사각형의 대칭 방향으로 정한다.
\begin{equation}
 \theta^*=\arg\min_{\theta\in\{\theta_b+k\cdot90^\circ\mid k\in\mathbb{Z}\}}
 \big|\operatorname{wrap}_{[-180^\circ,180^\circ)}(\theta-\theta_t)\big|.
\end{equation}
이 선택은 작업 영역의 접선과 블록 면을 함께 고려한다. 손목 회전의 안전성을 위해 실제 IK 해의 도달성도 확인하며, 과열 대응 전후의 온도 상승량과 연속 운용 시간은 별도로 비교하지 않았다.

\subsection{위치 보정에서 회전 재시도로의 변경}
초기에는 좌표 편차를 흡수하려고 반경·접선 방향의 경험적 보정과 네 위치 후보 $(0,15)$, $(-15,0)$, $(0,-15)$, $(15,0)$\,mm를 사용하였다. 후보 변경과 대상 재선택으로 동작을 이어 갔으나, 빈손 파지가 반복되거나 한 블록에 긴 시간이 소요되어 캘리브레이션과 파지 방향을 함께 점검하였다.

최종 기본 설정은 반경·전방 및 근거리 추가 보정을 0\,mm로 두고, 그리퍼 로컬 접선 방향의 5\,mm 보정을 사용한다. 위치 재시도 목록은 비워 두었다. 최초 실패나 하강 중단 뒤에는 같은 XY에서 파지 방향을 90도 회전한 후보를 한 번 시도한다. 두 회전 부호의 IK 해 가운데 도달 가능성과 손목 중립에 대한 편차를 고려해 하나를 선택한다. 회전 후보도 실패하면 재선택 또는 순회 보류로 이어진다.

\begin{figure}[htbp]
\centering
\includegraphics[width=0.7\linewidth]{camera_overlay}
\caption{관제 화면의 파지 방향 표시 예. 최종 방향 계산은 블록의 회전 사각형 방향과 베이스 중심 작업 영역의 접선을 이용하며, 화면의 선은 실제 턱 위치를 검출한 결과가 아니다.}
\label{fig:camera-overlay}
\end{figure}

관제 오버레이와 런타임 IK는 같은 정사각형 대칭 방향 선택 함수를 사용한다. 기본 파지 경로는 최초 관측으로 만든 후보를 실행하며, 파지 직전 실제 턱 위치를 재검출하여 XY를 보정하는 폐루프는 포함하지 않는다.

\subsection{도달성 검사와 접근 높이}
IK 결과는 모델 좌표에서의 위치 오차와 기울기 오차를 확인한 후 사용한다. 실행 전 계획의 허용 오차는 위치 15\,mm, 기울기 $6^\circ$로 설정하였다. 높은 접근점이 관절 한계에 걸릴 수 있으므로 최소 높이와 탐색 간격을 두고 가능한 접근 높이를 찾는다.

위치 15\,mm는 모델 안에서 계획 후보를 거르는 허용값이며, 실제 파지 여유를 측정한 값은 아니다. 그리퍼 개방 명령도 정규화 위치값이어서 실제 턱 간격으로 바로 환산할 수 없다. 따라서 IK 결과를 통과한 뒤에도 센서로 파지 상태를 확인한다.

인식 경로는 반경 320\,mm의 작업 영역 필터로 제어 대상을 제한한다. 이후 각 목표의 높이와 방향에 대해 IK를 풀어 실행 가능성을 판단한다. 관측 범위와 자세별 도달성을 두 단계로 검사하는 구조이다.

\subsection{측정 자세에서 시작하는 관절 보간}
웨이포인트 사이의 동작은 관절 공간에서 보간한다. 시작 자세 $\boldsymbol{q}_0$는 직전에 보낸 명령이 아니라 현재 측정값이며, 목표를 $\boldsymbol{q}_1$이라 하면 보간점은 다음과 같이 표현할 수 있다.
\begin{equation}
\boldsymbol{q}(\alpha)=(1-\alpha)\boldsymbol{q}_0+\alpha\boldsymbol{q}_1,
\qquad 0\leq\alpha\leq1.
\end{equation}
보간 과정에서 제어 주기(틱)당 관절 변화량을 제한하고, 로봇 입출력 단계의 상대 목표 클램프를 추가 적용한다. 측정 자세에서 시작하면 이전 동작이 목표에 충분히 도착하지 못한 경우에도 다음 궤적이 실제 자세와 다른 위치에서 시작해 급격한 명령 변화가 발생하는 것을 줄일 수 있다.

\begin{table}[htbp]
\centering
\caption{관절 제어와 파지 확인의 주요 설정값}
\begin{tabularx}{\linewidth}{@{}L{0.38\linewidth}L{0.19\linewidth}Y@{}}
\toprule
항목 & 설정값 & 의미 \\
\midrule
궤적 갱신 주파수 & 45\,Hz & 명령 생성의 목표 주파수 \\
틱당 일반 관절 변화 상한 & $2^\circ$ & 보간 단계의 제한 \\
하강 틱 변화 상한 & $0.6^\circ$ & 느린 접근을 위한 제한 \\
상대 목표 클램프 & 10.0 & 장비 입출력 계층의 추가 제한 \\
그리퍼 개방/닫힘 목표 & 95.0 / 2.0 & 정규화 위치 명령 \\
빈손 위치 경계 & 12.0 & 현재 파지 판정 설정 \\
파지 최소 부하 & 200.0 & 서보 부하값 기준 \\
파지 표본 수 / 간격 & 5회 / 0.05초 & 위치·부하 평균 계산 \\
\bottomrule
\end{tabularx}
\end{table}

명령 생성 목표 주파수는 30\,Hz에서 45\,Hz로 바꾸었다. 초기 180초 제한 시험에는 변경 전 설정을 사용하였다.

\section{상태 제어와 미션별 배치 전략}
\subsection{후보 실행과 파지 검증}
PICK에서는 기본 후보와 회전 재시도 후보의 IK 도달성을 검사한 뒤 실행 가능한 후보를 시도한다. 하강이 막히면 그리퍼를 열린 상태로 유지한 채 상공으로 후퇴하고, 정상 하강 후 빈손으로 닫히면 그리퍼를 다시 열고 후퇴하여 다음 후보를 시도한다. 후보를 소진하거나 개별 동작의 시간 제한에 걸리면 재선택 또는 순회 보류로 이어진다. PICK 내부의 후보 변경과 FSM의 대상 재선택은 서로 다른 복구 단계이다(그림~\ref{fig:pick-check}).

\begin{figure}[htbp]
\centering
\begin{tikzpicture}[
box/.style={draw,rounded corners,text width=3.7cm,minimum height=10mm,align=center,font=\small},
decision/.style={draw,diamond,aspect=2.6,text width=2.35cm,align=center,inner sep=1pt,font=\small},
arr/.style={-{Latex},thick}]
\node[box] (plan) at (0,0) {기본/회전 후보 계산};
\node[decision] (reach) at (0,-1.6) {도달 가능한\\후보 있음?};
\node[box] (attempt) at (0,-3.3) {정상 하강 후\\그리퍼 닫힘};
\node[decision] (held) at (0,-5.1) {$\bar p>12$ 그리고\\$\overline{|L|}\ge200$?};
\node[box] (success) at (0,-6.9) {PICK 통과\\필요시 후퇴 후 VERIFY 재확인};
\node[box,text width=2.8cm] (fail) at (5.1,-1.6) {PICK 실패\\재선택 / 순회 보류};
\node[decision] (more) at (5.1,-5.1) {남은 도달\\후보 있음?};
\draw[arr] (plan)--(reach);
\draw[arr] (reach)--node[right,font=\scriptsize]{예}(attempt);
\draw[arr] (reach)--node[above,font=\scriptsize]{아니오}(fail);
\draw[arr] (attempt)--(held);
\draw[arr] (held)--node[right,font=\scriptsize]{예}(success);
\draw[arr] (held)--node[above,font=\scriptsize]{아니오}(more);
\draw[arr] (more.north)--node[right,font=\scriptsize]{아니오}(fail.south);
\draw[arr] (more.north west)|-node[pos=0.75,above,font=\scriptsize]{예: 다음 후보}(attempt.east);
\end{tikzpicture}
\caption{PICK의 후보 선택과 정상 하강 후 센서 판정. 하강이 막힌 후보는 그리퍼를 닫지 않고 후퇴하여 다음 후보로 넘어간다. $\bar p$와 $\overline{|L|}$은 위치와 절대 부하의 표본 평균이며, 표시한 임계값은 평가에 사용한 설정이다. 개별 동작 시간 초과는 PICK 실패로 처리하며, 외부 180초 종료와 구분한다.}
\label{fig:pick-check}
\end{figure}

파지 판정에는 다섯 번 측정한 그리퍼 위치와 절대 부하의 평균을 사용하며 두 조건을 모두 만족해야 한다. VERIFY에서도 같은 조건을 검사하여 빈손 상태가 운반으로 이어지지 않게 한다. 실제 파지와 빈손 닫힘을 비교한 관찰은 절~\ref{sec:physical-grasp}에 제시한다.

\FloatBarrier
\subsection{추가 관측: 손목 영상과 마커 위치 추정}
\label{sec:additional-observation}
손목 카메라는 관제 영상과 데이터 수집에 활용하였다. 추가 관측 기능인 색상 게이트는 손목 영상에서 대상 색상이 검출되지 않으면 파지를 시작하지 않도록 하며, 미션 평가의 실제 제어에 적용하였다.

\begin{figure}[htbp]
\centering
\begin{subfigure}{0.48\linewidth}
\includegraphics[width=\linewidth]{wrist_camera_mount}
\caption{그리퍼에 재장착한 손목 카메라}
\end{subfigure}\hfill
\begin{subfigure}{0.48\linewidth}
\includegraphics[width=\linewidth]{wrist_camera_view}
\caption{관제 화면의 어깨(위)·손목(아래) 카메라 동시 표시}
\end{subfigure}
\caption{손목 카메라 장착과 관제 화면. 손목 영상은 관제·데이터 수집에 사용하며, 추가 색상 게이트의 입력으로도 활용하였다.}
\label{fig:wrist-camera}
\end{figure}

그리퍼에는 ArUco 마커를 부착하여 위치 추정을 실제 제어 실험에 적용하였다. 마커 기반 위치 추정은 상단 영상의 블록 검출과 함께 사용하는 추가 관측 기능이다. 추정 좌표를 자동 수집 데이터의 필드로 저장하고 활용하는 작업은 후속 과제로 남았다.

\begin{figure}[htbp]
\centering
\includegraphics[width=0.62\linewidth]{aruco_markers}
\caption{그리퍼에 부착한 ArUco 마커(35$\times$35\,mm)와 마커 시트, 체커보드, 대상 블록. 상단·하단 두 지점에 마커를 부착하여 그리퍼 자세를 함께 추정할 수 있도록 하였다.}
\label{fig:aruco}
\end{figure}
\FloatBarrier

\subsection{1차 미션: 지정 영역의 동적 슬롯 배치}
1차 미션은 지정 영역 내부에 먼 행 세 개와 가까운 행 두 개의 슬롯을 구성한다. 영역 네 꼭짓점과 정규화 좌표를 이용하여 슬롯 중심을 계산하고, 각 슬롯의 해제점과 상공 자세를 IK로 구한다. 기본 정규화 좌표는 $(0.14,0.20)$, $(0.50,0.20)$, $(0.80,0.20)$, $(0.32,0.76)$, $(0.68,0.76)$이다. 선택한 색은 예약한 슬롯을 유지하여 재시도 때 배치 위치가 불필요하게 바뀌는 것을 줄인다.

\begin{figure}[htbp]
\centering
\begin{tikzpicture}[x=0.055cm,y=0.055cm]
\draw[thick] (0,0) rectangle (200,100);
\draw[<->] (0,-12)--node[below]{200\,mm}(200,-12);
\draw[<->] (-12,0)--node[left,rotate=90,anchor=south]{100\,mm}(-12,100);
\foreach \x/\y/\name in {28/80/0,100/80/1,160/80/2,64/24/3,136/24/4}{
\draw[fill=black!6] (\x-20,\y-20) rectangle (\x+20,\y+20);
\fill (\x,\y) circle (1.5);
\node[above,font=\small] at (\x,\y+3) {\name};}
\node at (100,115) {먼 행};
\node at (100,-30) {로봇 쪽 / 가까운 행};
\end{tikzpicture}
\caption{200\,$\times$\,100\,mm 지정 영역의 3+2 슬롯 설계. 사각형은 40\,mm 블록의 명목 크기, 숫자는 슬롯 번호이다.}
\label{fig:slots}
\end{figure}

운반 경로는 파지 위치와 배치 위치의 높이·도달성을 고려하여 계산한다. 필요하면 베이스 쪽으로 당긴 경유점에서 더 높은 통과 자세를 만든 후 목표 슬롯 상공으로 이동한다. 해제 높이는 캘리브레이션의 평균 파지 높이에 가정한 여유 5\,mm를 더해 계산한다.

1차 미션은 HOME에서 새 프레임을 받아 지정 영역 밖의 검출이 연속 5초 동안 없으면 완료 상태로 전이한다. 중복되거나 오래된 프레임은 완료 시간에 누적하지 않는다. 이 판정은 관측된 블록 중심에 의존하므로, 검출 누락과 실제 경계 걸침을 포함한 물리적 완료 판정에는 한계가 있다.

\subsection{재시도와 실행 종료}
대상별 재시도 횟수는 한 순회 안에서만 제한한다. 남은 색이 모두 보류되면 시도 기록을 유지한 채 새 순회를 시작한다. 1차 미션 실행기는 내부 FSM 시간 제한을 비활성화한 상태로 외부 타이머를 사용하였다. 따라서 순회가 반복될 때 늘어나는 누적 재시도 시간은 외부 제한으로 관리한다.

제한시간 시험에서는 외부 타이머가 180초에 종료 신호를 보낸다. 종료 신호를 받은 뒤 그리퍼 개방과 홈 복귀를 시도하므로 전체 모터 동작이 180초 안에 끝나는지는 별도로 판정해야 한다.

\subsection{2차 미션: 적재 계획과 해제 전략}
2차 미션은 1차 미션의 대상 선택·파지·확인 흐름을 공유하고 목적지를 한 적재 좌표로 바꾼다. 계획기는 등록 영역 안의 좌표와 캘리브레이션의 파지 높이, 명목 블록 높이로 층별 접근·해제 자세를 IK로 계산한다. 평가에서는 접촉 하강을 비활성화하고, 계획한 해제 자세에 도달하면 그리퍼를 열어 다음 순회로 넘어갔다.

초기에는 관절 키프레임과 부하 변화로 접촉을 감지하는 하강 전략을 구성했으나, 최종 적재 시험에는 계산한 자세에서 직접 해제하는 방식을 사용하였다. 명목 높이와 배치 횟수로 만든 계획은 실제 타워 높이의 측정값이 아니므로, 해제 후 블록의 안착과 유지 상태를 별도로 관찰해야 한다.

적재 평가는 5단을 목표로 한 20회의 시험에서 수행 중 도달한 최대 높이를 집계하였다. 결과는 절~\ref{sec:stack-results}에 제시한다.

\subsection{Task3: 데이터셋 자동 수집 구조}
Task3은 CV+IK의 실행 중 관측과 동작을 기록하여 리더암 조작 없이 학습용 에피소드를 수집한다. 센서값과 검출 결과로 대상 선택·파지·배치를 판정하는 기존 실행 루프에 기록 기능을 연결하였다.

로봇 I/O 경계에 \texttt{RecordingRobotIO} 데코레이터를 두었다. 명령 직전의 관측과 안전 제한을 거쳐 실제 전송된 액션을 한 프레임으로 묶어 LeRobotDataset의 관절 상태, 설정한 카메라 영상, 액션과 작업 지시 형식으로 저장한다. PICK/VERIFY/TRANSPORT 코드는 1차 미션과 같다. 학습 자료의 행동값이 로봇 명령과 일치하도록 안전 제한 전의 요청값 대신 실제 전송값을 기록한다.

Task3의 에피소드는 대상 선택을 마친 HOME 자세에서 시작하여 파지$\to$운반$\to$배치$\to$HOME 복귀로 끝난다. 대상 선택을 위한 카메라 대기는 기록 구간에서 제외하며, 센서 파지 확인을 통과하고 배치와 복귀까지 끝난 사이클을 저장한다. 개별 에피소드의 실제 안착 여부를 독립 영상 판정한 라벨과는 구분한다. 1차/2차 미션과 달리 회전·위치 후보로 재시도하지 않고 최초 파지 1회만 시도하여, 복구 동작을 제외한 실행을 수집하도록 설계하였다. 지정 영역 밖의 블록이 5초간 탐지되지 않으면 한 라운드가 끝난 것으로 보고 작업자에게 블록 재배치를 요청한다.

저장 자료는 관절·행동 데이터, 카메라 영상과 에피소드 메타데이터로 구성된다(표~\ref{tab:task3-structure}).

\begin{table}[htbp]
\centering
\caption{자동 수집 데이터의 구성}\label{tab:task3-structure}
\begin{tabularx}{\linewidth}{@{}L{0.22\linewidth}L{0.20\linewidth}Y@{}}
\toprule 자료 & 형식 & 저장 내용 \\
\midrule
관절·행동 & parquet & 관측한 관절 상태와 실제 전송한 액션 \\
카메라 영상 & MP4 & 설정한 카메라의 관측 영상 \\
메타데이터 & JSON, parquet & 에피소드 구분, 작업 지시와 통계 \\
\bottomrule
\end{tabularx}
\end{table}

\subsection{안전 처리와 실행 기록}
안전장치는 관절 보간의 변화량 제한, 상대 목표 클램프, IK 게이트와 파지 확인에 나누어 적용하였다. 하강 중 턱이 물체에 걸렸는데도 명령을 계속 보내면 블록을 밀거나 팔이 구속될 수 있다. 측정 관절과 명령의 차이를 검사해 하강을 중단하며, PICK에서는 그리퍼를 닫지 않고 상공으로 후퇴한 뒤 다음 후보를 시도한다. 하강 중 명령 추종 오차의 상한은 8.0, 하강 종료 시 목표 미도달 판정 경계는 4.0으로 조정하였다. 두 값은 로봇 입출력의 관절 명령 단위를 따르는 설정값이다.

정상 종료와 예외 발생 시 모두 그리퍼 개방, 홈 복귀와 연결 해제를 시도하도록 종료 절차를 구성하였다. 통신이나 장비 오류가 지속되면 정리 동작이 실패할 수 있으므로 오류와 최종 상태를 함께 기록한다. 연결 해제 후에는 팔의 자세 유지를 위해 토크를 유지한다. 모터를 구동하지 않는 사전 경로 검사와 실제 장비 시험은 별도로 수행하였다.

실험마다 적용 설정, 대응점·관절값, 상태 전이와 단계별 사진을 저장하여 실패 단계와 실행 조건을 추적하였다.

\Needspace{27\baselineskip}
\section{개발 일정과 구성원별 수행 내용}
\subsection{개발 경과}
개발 경과와 산출물은 표~\ref{tab:timeline}과 같다.
\begin{longtable}{@{}L{0.20\linewidth}L{0.70\linewidth}@{}}
\caption{주요 개발 경과와 산출물}\label{tab:timeline}\\
\toprule 시기 & 수행 내용 및 산출물 \\\midrule
\endfirsthead
\toprule 시기 & 수행 내용 및 산출물 \\\midrule
\endhead
3--5월 & 초기 물체 수집 주제 검토, 시스템 구상과 착수보고서 작성 \\
6월 & SO-101 과제로 변경, 장비 수령, 리더·팔로워 환경과 데이터 수집 기반 구성 \\
7월 & SmolVLA·ACT 학습·실행, 원격 추론과 카메라 도구 개발, 중간보고서 작성 \\
8월 초 & 산업체 서면 자문, 구현·검증 구분과 접근법 선정·비용 비교 의견 수렴 \\
8월 말 & CV+IK 파지 경로, 캘리브레이션·기구학 도구 개발, 손목 방향과 위치 오차 분석 \\
9월 / 통합 시연 & 검출과 3+2 슬롯 배치 기능 통합, 1차 미션 시연 \\
9월 / 실기 평가 & 카메라 재고정, 9점 캘리브레이션, 제한시간·단일 블록 시험, 실험 기록 정리 \\
9월 / 결과 정리 & 최종 시스템 결과 분석, 최종보고서 작성과 자문의견 반영 \\
\bottomrule
\end{longtable}

\Needspace{31\baselineskip}
\subsection{역할 분담과 공동 작업}
비전과 미션 제어, 관제·캘리브레이션 도구, 평가와 기술 문서를 나누어 담당했으며 장비 통합과 결과 검토는 공동으로 진행하였다(표~\ref{tab:roles}).

\begin{longtable}{@{}L{0.14\linewidth}L{0.82\linewidth}@{}}
\caption{구성원별 수행 내용}\label{tab:roles}\\
\toprule 구성원 & 수행 내용 \\\midrule
\endfirsthead
\toprule 구성원 & 수행 내용 \\\midrule
\endhead
윤민석 & 블록별 채도·밝기와 영역 필터 조정(도달 범위 밖·지정 영역 내 블록 제외). 파지·적재 위치의 mm 단위 보정과 블록 각도 기반 그리퍼 회전 정렬.\newline
손목 서보 과열 분석과 제어 개선. 파지 실패 재시도, 성공 후 운반과 미션 종료 조건 구현.\newline
1차 미션 제어 경로 통합 및 졸업과제 포스터 제작. \\
이동근 & FSM 파지 조건과 도달 실패 시 대체 후보 처리 보완. 카메라 웹 뷰어 및 검출·파지 후보·FK 위치의 실시간 시각화.\newline
대응점·관절값 기록과 수집 재개를 지원하는 캘리브레이션 도구 개발. 영상 변화 감지에 따른 실험 이미지 자동 저장 구현.\newline
LeRobot 의존성과 uv 환경 통합, Python CLI 실행 도구 정리. 설정·로그·캘리브레이션 자료 관리 및 중간·최종보고서 작성. \\
김주환 & SmolVLA·ACT의 한계 분석과 규칙 기반 FSM 전환 근거 정리. 정량 평가 체계 및 실행 기록 양식 작성·배포.\newline
캘리브레이션 시행착오 문서화·공유. 실행 환경 통합 시험과 미션 실행 검증.\newline
코드 검토, 개발 규칙·변경 이력 관리 및 최종보고서·최종 시연 PPT 작성. \\
공동 수행 & 실제 장비 통합, 데이터 수집과 실패 관찰. 접근법 논의, 결과 검토와 보고 자료 정리. \\
\bottomrule
\end{longtable}
\FloatBarrier
